\ifdefined\gkver\else\def\gkver{student}\fi
\documentclass[\gkver]{gaokaozhenti}
\begin{document}

\chapter{2008年山东理综}
%% paperType: 理综物理部分

\begin{enumerate}
\item
%% number: 16
%% paperName: 2008年山东理综第 16 题 4 分
%% typeId: 1
%% type: 单选题
%% chapter: 力物体的平衡
%% point: 三力平衡
%% method: 无
%% score: 4
%% degree: 650
%% duplicateId: 0
%% body:
用轻弹簧竖直悬挂的质量为 $m$ 的物体，静止时弹簧伸长量为 $L_{0}$。现用该弹簧沿斜面方向拉住质量为 $2m$ 的物体，系统静止时弹簧伸长量也为 $L_{0}$，斜面倾角为 $30^{\circ}$，如图所示。则物体所受摩擦力\xzanswer{A}

\onepicture[w=3.5cm]{figs/16.png}

\fourchoices[answer=A]
{等于零}
{大小为 $\frac{1}{2}mg$，方向沿斜面向下}
{大小为 $\frac{\sqrt{3}}{2}mg$，方向沿斜面向上}
{大小为 $mg$，方向沿斜面向上}

%% answer: A
%% memo:
\memoanswer{竖直挂时 $mg=k\Delta x$，当质量为 $2m$ 放到斜面上时，$2mg\sin30^{\circ}=f+k\Delta x$，因两次弹簧长度一样，所以 $\Delta x$ 也一样。解这两个方程可得，物体受到的摩擦力为零，A 正确。}

\item
%% number: 17
%% paperName: 2008年山东理综第 17 题 4 分
%% typeId: 2
%% type: 多选题
%% chapter: 直线运动
%% point: 运动图象
%% method: 无
%% score: 4
%% degree: 650
%% duplicateId: 0
%% body:
质量为 $1500\Ukg$ 的汽车在平直的公路上运动，$v-t$ 图象如图所示。由此可求\xzanswer{ABD}

\onepicture[w=4.5cm]{figs/17.png}

\fourchoices[answer=ABD]
{前 $25\Us$ 内汽车的平均速度}
{前 $10\Us$ 内汽车的加速度}
{前 $10\Us$ 内汽车所受的阻力}
{$15\sim25\Us$ 内合外力对汽车所做的功}

%% answer: ABD
%% memo:
\memoanswer{通过 $v-t$ 图象的面积就是物体的位移，所以能求出位移，还知道时间，所以能求出平均速度，A 对。$v-t$ 图象的斜率就是物体的加速度，所以能得到 $10\Us$ 内的加速度，B 对。不知道汽车的牵引力，所以得不出受到的阻力，C 错。$15\Us$ 到 $25\Us$ 汽车的初速度和末速度均已知，由动能定理，可以得出合外力做的功，D 对。}

\item
%% number: 18
%% paperName: 2008年山东理综第 18 题 4 分
%% typeId: 2
%% type: 多选题
%% chapter: 曲线运动
%% point: 人造地球卫星
%% method: 无
%% score: 4
%% degree: 650
%% duplicateId: 0
%% body:
据报道，我国数据中继卫星“天链一号01星”于 2008 年 4 月 25 日在西昌卫星发射中心发射升空，经过 4 次变轨控制后，于 5 月 1 日成功定点在东经 $77^{\circ}$ 赤道上空的同步轨道。关于成功定点后的“天链一号01星”，下列说法正确的是\xzanswer{BC}

\fourchoices[answer=BC]
{运行速度大于 $7.9\Ukms$}
{离地面高度一定，相对地面静止}
{绕地球运行的角速度比月球绕地球运行的角速度大}
{向心加速度与静止在赤道上物体的向心加速度大小相等}

%% answer: BC
%% memo:
\memoanswer{由题目可以得出“天链一号卫星”是地球同步卫星，运行速度要小于 $7.9\Ukms$，而它的位置在赤道上空，高度一定，A 错 B 对。由 $\omega=\frac{2\pi}{T}$ 可知，C 对。由 $a=\frac{GM}{R^{2}}$ 可知，D 错。}

\item
%% number: 19
%% paperName: 2008年山东理综第 19 题 4 分
%% typeId: 1
%% type: 单选题
%% chapter: 运动和力
%% point: 变加速直线运动
%% method: 无
%% score: 4
%% degree: 450
%% duplicateId: 0
%% body:
直升机悬停在空中向地面投放装有救灾物资的箱子，如图所示。设投放初速度为零，箱子所受的空气阻力与箱子下落速度的平方成正比，且运动过程中箱子始终保持图示姿态。在箱子下落过程中，下列说法正确的是\xzanswer{C}

\onepicture[w=2.7cm]{figs/19.png}

\fourchoices[answer=C]
{箱内物体对箱子底部始终没有压力}
{箱子刚从飞机上投下时，箱内物体受到的支持力最大}
{箱子接近地面时，箱内物体受到的支持力比刚投下时大}
{若下落距离足够长，箱内物体有可能不受底部支持力而“飘起来”}

%% answer: C
%% memo:
\memoanswer{因为受到阻力，箱子不是处于完全失重状态，所以箱内物体对支持面有压力，A 错。由于箱子阻力和下落的速度成二次方关系，最终将匀速运动，此时受到的支持力等于重力，B、D 错，C 对。}

\item
%% number: 20
%% paperName: 2008年山东理综第 20 题 4 分
%% typeId: 1
%% type: 单选题
%% chapter: 交流电
%% point: 交流电
%% method: 无
%% score: 4
%% degree: 650
%% duplicateId: 0
%% body:
图 1、图 2 分别表示两种电压的波形，其中图 1 所示电压按正弦规律变化。下列说法正确的是\xzanswer{C}

\onepicture[w=10.0cm]{figs/20.png}

\fourchoices[answer=C]
{图 1 表示交流电，图 2 表示直流电}
{两种电压的有效值相等}
{图 1 所示电压的瞬时值表达式为 $u=311\sin100\pi t\UV$}
{图 1 所示电压经匝数比为 $10:1$ 的变压器变压后，频率变为原来的 $\frac{1}{10}$}

%% answer: C
%% memo:
\memoanswer{交流电的大小和方向都随时间变化，在 $t$ 轴的上方为正，下方为负，A 错。有效值 $E=\frac{E_{m}}{\sqrt{2}}$ 只对正弦交流电适用，两种电压的最大值一样，所以 B 错。由图可知，C 对。变压之后频率不变，D 错。}

\item
%% number: 21
%% paperName: 2008年山东理综第 21 题 4 分
%% typeId: 2
%% type: 多选题
%% chapter: 电场
%% point: 电势能电势差电场力做功
%% method: 无
%% score: 4
%% degree: 450
%% duplicateId: 0
%% body:
如图所示，在 $y$ 轴上关于 $O$ 点对称的 $A$、$B$ 两点有等量同种点电荷 $+Q$，在 $x$ 轴上 $C$ 点有点电荷 $-Q$，且 $CO=OD$，$\angle ADO=60^{\circ}$。下列判断正确的是\xzanswer{BD}

\onepicture[w=2.7cm]{figs/21.png}

\fourchoices[answer=BD]
{$O$ 点电场强度为零}
{$D$ 点电场强度为零}
{若将点电荷 $+q$ 从 $O$ 移向 $C$，电势能增大}
{若将点电荷 $-q$ 从 $O$ 移向 $C$，电势能增大}

%% answer: BD
%% memo:
\memoanswer{电场是矢量，叠加遵循平行四边形定则，由 $E=\frac{kQ}{r^{2}}$ 和几何关系可以得出，$O$ 点电场强度不为零，$D$ 点电场强度为零，A 错 B 对。在 $O\to C$ 之间，合场强的方向向左，把负电荷从 $O$ 移动到 $C$，电场力做负功，电势能增加，C 错 D 对。}

\item
%% number: 22
%% paperName: 2008年山东理综第 22 题 4 分
%% typeId: 1
%% type: 单选题
%% chapter: 电磁感应
%% point: 电磁感应能量分析
%% method: 无
%% score: 4
%% degree: 450
%% duplicateId: 0
%% body:
两根足够长的光滑导轨竖直放置，间距为 $L$，底端接阻值为 $R$ 的电阻。将质量为 $m$ 的金属棒悬挂在一个固定的轻弹簧下端，金属棒和导轨接触良好，导轨所在平面与磁感应强度为 $B$ 的匀强磁场垂直，如图所示。除电阻 $R$ 外其余电阻不计。现将金属棒从弹簧原长位置由静止释放，则\xzanswer{C}

\onepicture[w=3.0cm]{figs/22.png}

\fourchoices[answer=C]
{释放瞬间金属棒的加速度等于重力加速度 $g$}
{金属棒向下运动时，流过电阻 $R$ 的电流方向为 $a\to b$}
{金属棒的速度为 $v$ 时，所受的安培力大小为 $F=\frac{B^{2}L^{2}v}{R}$}
{电阻 $R$ 上产生的总热量等于金属棒重力势能的减少}

%% answer: C
%% memo:
\memoanswer{在释放的瞬间，速度为零，不受安培力的作用，只受到重力，A 对。由右手定则可得，电流的方向从 $b$ 到 $a$，B 错。当速度为 $v$ 时，产生的电动势为 $E=Blv$，受到的安培力为 $F=BIL$，计算可得 $F=\frac{B^{2}L^{2}v}{R}$，C 对。在运动的过程中，是弹簧的弹性势能、重力势能和内能的转化，D 错。（本题原卷参考答案页标 AC，而题后答案及平台数据标单选 C，二者不一致；答案命令按平台数据取 C，解析照原卷如实保留。）}

\item
%% number: 23
%% paperName: 2008年山东理综第 23 题 12 分
%% typeId: 4
%% type: 实验题
%% chapter: 电路
%% point: 电路实验
%% method: 无
%% score: 12
%% degree: 450
%% duplicateId: 0
%% body:
2007 年诺贝尔物理学奖授予了两位发现“巨磁电阻”效应的物理学家。材料的电阻随磁场的增加而增大的现象称为磁阻效应，利用这种效应可以测量磁感应强度。

若图 1 为某磁敏电阻在室温下的电阻-磁感应强度特性曲线，其中 $R_{B}$、$R_{0}$ 分别表示有、无磁敏电阻的阻值。为了测量磁感应强度 $B$，需先测量磁敏电阻处于磁场中的电阻值 $R_{B}$。请按要求完成下列实验。

\onepicture[w=4.0cm]{figs/23a.png}

\begin{enumerate}
	\item
	设计一个可以测量磁场中该磁敏电阻阻值的电路，在图 2 的虚线框内画出实验电路原理图（磁敏电阻及所处磁场已给出，待测磁场磁感应强度大小约为 $0.6\sim1.0\UT$，不考虑磁场对电路其它部分的影响）。要求误差较小。
	
	\onepicture[w=3.2cm]{figs/23b.png}
	
	提供的器材如下：
	
	A．磁敏电阻，无磁场时阻值 $R_{0}=150\UO$
	
	B．滑动变阻器 $R$，全电阻约 $20\UO$
	
	C．电流表，量程 $2.5\UmA$，内阻约 $30\UO$
	
	D．电压表，量程 $3\UV$，内阻约 $3\UkO$
	
	E．直流电源 $E$，电动势 $3\UV$，内阻不计
	
	F．开关 $S$，导线若干
	
	\item
	正确接线后，将磁敏电阻置入待测磁场中，测量数据如下表：
	
	\begin{table}[htbp]
	\centering
	\begin{tblr}{|c|c|c|c|c|c|c|}
	\hline
	 & 1 & 2 & 3 & 4 & 5 & 6 \\
	\hline
	$U$（$\UV$） & 0.00 & 0.45 & 0.91 & 1.50 & 1.79 & 2.71 \\
	\hline
	$I$（$\UmA$） & 0.00 & 0.30 & 0.60 & 1.00 & 1.20 & 1.80 \\
	\hline
	\end{tblr}
	\end{table}
	
	根据上表可求出磁敏电阻的测量值 $R_{B}=$ \tkanswer[1.5cm]{1500} $\UO$，结合图 1 可知待测磁场的磁感应强度 $B=$ \tkanswer[1.5cm]{0.90} $\UT$。
	
	\item
	试结合图 1 简要回答，磁感应强度 $B$ 在 $0\sim0.2\UT$ 和 $0.4\sim1.0\UT$ 范围内磁敏电阻阻值的变化规律有何不同？
	
	\item
	某同学查阅相关资料时看到了图 3 所示的磁敏电阻在一定温度下的电阻-磁感应强度特性曲线（关于纵轴对称），由图线可以得到什么结论？
	
	\onepicture[w=4.0cm]{figs/23c.png}
\end{enumerate}

%% answer:
\jdanswer{
\begin{enumerate}
	\item
	电路图如图所示。
	
	\item
	$1500$，$0.90$
	
	\item
	在 $0\sim0.2\UT$ 范围内，磁敏电阻的阻值随磁感应强度非线性变化（或不均匀变化）；在 $0.4\sim1.0\UT$ 范围内，磁敏电阻的阻值随磁感应强度线性变化（或均匀变化）。
	
	\item
	磁场反向，磁敏电阻的阻值不变。
\end{enumerate}
}

%% memo:
\memoanswer{
（1）实验电路如图所示。

\onepicture[w=5.0cm]{figs/23_an.jpg}

（2）根据表中数据可求出磁敏电阻的测量值 $R_{B}=1500\UO$，结合图 1 可知待测磁场的磁感应强度 $B=0.90\UT$。

（3）在 $0\sim0.2\UT$ 范围内，磁敏电阻的阻值随磁感应强度非线性变化（或不均匀变化）；在 $0.4\sim1.0\UT$ 范围内，磁敏电阻的阻值随磁感应强度线性变化（或均匀变化）。

（4）磁场反向，磁敏电阻的阻值不变。
}

\item
%% number: 24
%% paperName: 2008年山东理综第 24 题 15 分
%% typeId: 6
%% type: 计算题
%% chapter: 曲线运动
%% point: 竖直平面圆周运动
%% method: 无
%% score: 15
%% degree: 650
%% duplicateId: 0
%% body:
某兴趣小组设计了如图所示的玩具轨道，其中“2008”四个等高数字用内壁光滑的薄壁细圆管弯成，固定在竖直平面内（所有数字均由圆或半圆组成，圆半径比细管的内径大得多），底端与水平地面相切。弹射装置将一个小物体（可视为质点）以 $v_{a}=5\Ums$ 的水平初速度由 $a$ 点弹出，从 $b$ 点进入轨道，依次经过“8002”后从 $p$ 点水平抛出。小物体与地面 $ab$ 段间的动摩擦因数 $\mu=0.3$，不计其它机械能损失。已知 $ab$ 段长 $L=1.5\Um$，数字“0”的半径 $R=0.2\Um$，小物体质量 $m=0.01\Ukg$，$g=10\Umsq$。求：

\begin{enumerate}
	\item
	小物体从 $p$ 点抛出后的水平射程。
	\item
	小物体经过数字“0”的最高点时管道对小物体作用力的大小和方向。
\end{enumerate}

\onepicture[w=8.0cm, align=right]{figs/24.png}

%% answer:
\jdanswer{
\begin{enumerate}
	\item
	$s=0.8\Um$
	
	\item
	$F=0.3\UN$，方向竖直向下
\end{enumerate}
}

%% memo:
\memoanswer{
（1）设小物体运动到 $p$ 点时的速度大小为 $v$，对小物体由 $a$ 到 $p$ 过程应用动能定理得
\[ -\mu mgL-2mgR=\frac{1}{2}mv^{2}-\frac{1}{2}mv_{a}^{2}, \]
\[ 2R=\frac{1}{2}gt^{2},\qquad s=vt. \]
联立以上各式，代入数据解得 $s=0.8\Um$。

（2）设在数字“0”的最高点时管道对小物体的作用力大小为 $F$，取竖直向下为正方向，有
\[ F+mg=\frac{mv^{2}}{R}, \]
代入数据解得 $F=0.3\UN$，方向竖直向下。
}

\item
%% number: 25
%% paperName: 2008年山东理综第 25 题 18 分
%% typeId: 6
%% type: 计算题
%% chapter: 磁场
%% point: 带电粒子在复合场中的运动
%% method: 无
%% score: 18
%% degree: 450
%% duplicateId: 0
%% body:
两块足够大的平行金属极板水平放置，极板间加有空间分布均匀、大小随时间周期性变化的电场和磁场，变化规律分别如图 1、图 2 所示（规定垂直纸面向里为磁感应强度的正方向）。在 $t=0$ 时刻由负极板释放一个初速度为零的带负电的粒子（不计重力）。若电场强度 $E_{0}$、磁感应强度 $B_{0}$、粒子的比荷 $\frac{q}{m}$ 均已知，且 $t_{0}=\frac{2\pi m}{qB_{0}}$，两板间距 $h=\frac{10\pi^{2}mE_{0}}{qB_{0}^{2}}$。

\begin{enumerate}
	\item
	求粒子在 $0\sim t_{0}$ 时间内的位移大小与极板间距 $h$ 的比值；
	\item
	求粒子在极板间做圆周运动的最大半径（用 $h$ 表示）；
	\item
	若板间电场强度 $E$ 随时间的变化仍如图 1 所示，磁场的变化改为如图 3 所示，试画出粒子在板间运动的轨迹图（不必写计算过程）。
\end{enumerate}

\onepicture[w=9.0cm]{figs/25.png}

%% answer:
\jdanswer{
\begin{enumerate}
	\item
	$\frac{s_{1}}{h}=\frac{1}{5}$
	
	\item
	$R_{2}=\frac{2h}{5\pi}$
	
	\item
	轨迹图如图所示。
\end{enumerate}
}

%% memo:
\memoanswer{
（1）设粒子在 $0\sim t_{0}$ 时间内运动的位移大小为 $s_{1}$，则
\[ s_{1}=\frac{1}{2}at_{0}^{2},\qquad a=\frac{qE_{0}}{m}. \]
又 $t_{0}=\frac{2\pi m}{qB_{0}}$，$h=\frac{10\pi^{2}mE_{0}}{qB_{0}^{2}}$，解得 $\frac{s_{1}}{h}=\frac{1}{5}$。

（2）粒子在 $t_{0}\sim2t_{0}$ 时间内只受洛伦兹力作用，且速度与磁场方向垂直，所以粒子做匀速圆周运动。设运动速度大小为 $v_{1}$，轨道半径为 $R_{1}$，周期为 $T$，则
\[ v_{1}=at_{0},\qquad qv_{1}B_{0}=\frac{mv_{1}^{2}}{R_{1}}, \]
解得 $R_{1}=\frac{h}{5\pi}$。又 $T=\frac{2\pi m}{qB_{0}}=t_{0}$，即粒子在 $t_{0}\sim2t_{0}$ 时间内恰好完成一个周期的圆周运动。

在 $2t_{0}\sim3t_{0}$ 时间内，粒子做初速度为 $v_{1}$ 的匀加速直线运动，设位移大小为 $s_{2}$，则
\[ s_{2}=v_{1}t_{0}+\frac{1}{2}at_{0}^{2}, \]
解得 $s_{2}=\frac{3}{5}h$。

由于 $s_{1}+s_{2}<h$，所以粒子在 $3t_{0}\sim4t_{0}$ 时间内继续做匀速圆周运动，设速度大小为 $v_{2}$，半径为 $R_{2}$，则
\[ v_{2}=v_{1}+at_{0},\qquad qv_{2}B_{0}=\frac{mv_{2}^{2}}{R_{2}}, \]
解得 $R_{2}=\frac{2h}{5\pi}$。由于 $s_{1}+s_{2}+R_{2}<h$，粒子恰好又完成一个周期的圆周运动；在 $4t_{0}\sim5t_{0}$ 时间内，粒子运动到正极板。因此粒子运动的最大半径 $R_{2}=\frac{2h}{5\pi}$。

（3）粒子在板间运动的轨迹如图所示。

\onepicture[w=5.0cm]{figs/25_an.jpg}
}

\item
%% number: 36
%% paperName: 2008年山东理综第 36 题 8 分
%% typeId: 6
%% type: 计算题
%% chapter: 气体性质
%% point: 玻意耳定律
%% method: 无
%% score: 8
%% degree: 650
%% duplicateId: 0
%% body:
【物理——物理 3—3】喷雾器内有 $10\UL$ 水，上部封闭有 $1\Uatm$ 的空气 $2\UL$。关闭喷雾阀门，用打气筒向喷雾器内再充入 $1\Uatm$ 的空气 $3\UL$（设外界环境温度一定，空气可看作理想气体）。

\begin{enumerate}
	\item
	当水面上方气体温度与外界温度相等时，求气体压强，并从微观上解释气体压强变化的原因。
	\item
	打开喷雾阀门，喷雾过程中封闭气体可以看成等温膨胀，此过程气体是吸热还是放热？简要说明理由。
\end{enumerate}

\onepicture[w=5.0cm, align=right]{figs/36.png}

%% answer:
\jdanswer{
\begin{enumerate}
	\item
	$p_{2}=2.5\Uatm$。微观解释：温度不变，分子平均动能不变，单位体积内分子数增加，所以压强增加。
	
	\item
	吸热。气体对外做功而内能不变，根据热力学第一定律可知气体吸热。
\end{enumerate}
}

%% memo:
\memoanswer{
（1）设气体初态压强为 $p_{1}$，体积为 $V_{1}$；末态压强为 $p_{2}$，体积为 $V_{2}$，由玻意耳定律
\[ p_{1}V_{1}=p_{2}V_{2}, \]
代入数据得 $p_{2}=2.5\Uatm$。

微观解释：温度不变，分子平均动能不变，单位体积内分子数增加，所以压强增加。

（2）吸热。气体对外做功而内能不变，根据热力学第一定律可知气体吸热。
}

\item
%% number: 37
%% paperName: 2008年山东理综第 37 题 8 分
%% typeId: 6
%% type: 计算题
%% chapter: 光学
%% point: 光的折射
%% method: 无
%% score: 8
%% degree: 650
%% duplicateId: 0
%% body:
【物理——物理 3—4】麦克斯韦在 1865 年发表的《电磁场的动力学理论》一文中揭示了电、磁现象与光的内在联系及统一性，即光是电磁波。

\begin{enumerate}
	\item
	一单色光波在折射率为 1.5 的介质中传播，某时刻电场横波图象如图 1 所示，求该光波的频率。
	\item
	图 2 表示两面平行玻璃砖的截面图，一束平行于 $CD$ 边的单色光入射到 $AC$ 界面上，$a$、$b$ 是其中的两条平行光线。光线 $a$ 在玻璃砖中的光路已给出。画出光线 $b$ 从玻璃砖中首次出射的光路图，并标出出射光线与界面法线夹角的度数。
\end{enumerate}

\onepicture[w=6.0cm]{figs/37a.png}

\onepicture[w=6.0cm]{figs/37b.png}

%% answer:
\jdanswer{
\begin{enumerate}
	\item
	$f=5\times10^{14}\UHz$
	
	\item
	光路图如图所示。
\end{enumerate}
}

%% memo:
\memoanswer{
（1）设光在介质中的传播速度为 $v$，波长为 $\lambda$，频率为 $f$，则
\[ f=\frac{v}{\lambda},\qquad v=\frac{c}{n}. \]
联立以上两式得 $f=\frac{c}{n\lambda}$。从波形图上读出波长 $\lambda=4\times10^{-7}\Um$，代入数据解得
\[ f=5\times10^{14}\UHz. \]

（2）光路如图所示。

\onepicture[w=5.0cm]{figs/37_an.jpg}
}

\item
%% number: 38
%% paperName: 2008年山东理综第 38 题 8 分
%% typeId: 6
%% type: 计算题
%% chapter: 运动和力
%% point: 动量守恒定律
%% method: 无
%% score: 8
%% degree: 650
%% duplicateId: 0
%% body:
【物理——物理 3—5】

\begin{enumerate}
	\item
	在氢原子光谱中，电子从较高能级跃迁到 $n=2$ 能级发出的谱线属于巴耳末线系。若一群氢原子自发跃迁时发出的谱线中只有 2 条属于巴耳末线系，则这群氢原子自发跃迁时最多可发生 \tkanswer[1.5cm]{6} 条不同频率的谱线。
	\item
	一个物体静置于光滑水平面上，外面扣一质量为 $M$ 的盒子，如图 1 所示。现给盒子一初速度 $v_{0}$，此后，盒子运动的 $v-t$ 图象呈周期性变化，如图 2 所示。请据此求盒内物体的质量。
\end{enumerate}

\onepicture[w=9.0cm]{figs/38.png}

%% answer:
\jdanswer{
\begin{enumerate}
	\item
	$6$
	
	\item
	$m=M$
\end{enumerate}
}

%% memo:
\memoanswer{
（1）$6$。

（2）设物体的质量为 $m$，$t_{0}$ 时刻物体受盒子碰撞获得速度 $v$，根据动量守恒定律
\[ Mv_{0}=mv. \]
$3t_{0}$ 时刻物体与盒子右壁碰撞使盒子速度又变为 $v_{0}$，说明碰撞是弹性碰撞，则
\[ \frac{1}{2}Mv_{0}^{2}=\frac{1}{2}mv^{2}. \]
联立以上两式解得 $m=M$。（也可通过图象分析得出 $v_{0}=v$，结合动量守恒，得出正确结果。）
}

\end{enumerate}

\end{document}
